\section{Nhị thức Newtơn}

\subsection{Đề 1}
\Opensolutionfile{ans}[ans]
\caulc
\begin{ex}%[0D8N3-1]
	Trong khai triển nhị thức Newtơn của $(a+b)^4$ có bao nhiêu số hạng?
	\choice
	{$6$}
	{$3$}
	{\True $5$}
	{$4$}
	\loigiai{
		Trong khai triển nhị thức Newtơn của $(a+b)^4$ có $4+1=5$ số hạng.
	}
\end{ex}
\begin{ex}%[0D8H3-2]
	Trong khai triển nhị thức Newtơn của $(2x-3)^4$ có bao nhiêu số hạng?
	\choice
	{$6$}
	{$3$}
	{\True $5$}
	{$4$}
	\loigiai{
		Trong khai triển nhị thức Newtơn của $(2x-3)^4$ có $4+1=5$ số hạng.
	}
\end{ex}
\begin{ex}%[0D8N3-2]
	Trong khai triển nhị thức Newtơn của $(a+b)^4$, số hạng tổng quát của khai triển là
	\choice
	{$\mathrm{C}_4^{k-1}a^kb^{5-k}$}
	{\True $\mathrm{C}_4^ka^{4-k}b^k$}
	{$\mathrm{C}_4^{k+1}a^{5-k}b^{k+1}$}
	{$\mathrm{C}_4^ka^{4-k}b^{4-k}$}
	\loigiai{
		Số hạng tổng quát của khai triển $(a+b)^4$ là $\mathrm{C}_n^ka^{n-k}b^k=\mathrm{C}_4^ka^{4-k}b^k$.
	}
\end{ex}
\begin{ex}%[0D8H3-2]
	Trong khai triển nhị thức Newtơn của $(2x-3)^4$, số hạng tổng quát của khai triển là
	\choice
	{$\mathrm{C}_4^k2^k3^{4-k}\cdot x^{4-k}$}
	{\True $\mathrm{C}_4^k2^{4-k}(-3)^k\cdot x^{4-k}$}
	{$\mathrm{C}_4^k2^{4-k}3^k\cdot x^{4-k}$}
	{$\mathrm{C}_4^k2^k\left(-3\right)^{4-k}\cdot x^{4-k}$}
	\loigiai{
		Số hạng tổng quát của khai triển $(2x-3)^4$ là $\mathrm{C}_4^k(2y)^{4-k}\left(-3\right)^k=\mathrm{C}_4^k2^{4-k}(-3)^k\cdot x^{4-k}$.
	}
\end{ex}
\begin{ex}%[0D8H3-5]
	Tính tổng các hệ số trong khai triển nhị thức Newtơn của $(1-2x)^4$.
	\choice
	{\True $1$}
	{$-1$}
	{$81$}
	{$-81$}
	\loigiai{
		Khai triển $(2x-3)^4$ có dạng $(2x-3)^4=a_0x^4+a_1x^3+a_2x^2+a_3x+a_4$ \hfill $(1)$\\
		Từ $(1)$, cho $x=1$ ta được $a_0+a_1+a_2+a_3+a_4=(2\cdot 1-3)^4=1$.\\
		Vậy $S=1$.
	}
\end{ex}
\begin{ex}%[0D8H3-4]
	Trong khai triển nhị thức Newtơn của $(1+3x)^4$, số hạng thứ $2$ theo số mũ tăng dần của $x$ là
	\choice
	{$108x$}
	{$54x^2$}
	{$1$}
	{\True $12x$}
	\loigiai{
		Ta có $(1+3x)^4=\sum\limits_{k=0}^4\mathrm{C}_4^k(3x)^k=\sum\limits_{k=0}^4\mathrm{C}_4^k3^kx^k$.\\
		Do đó số hạng thứ $2$ theo số mũ tăng dần của $x$ ứng với $k=1$, tức là $\mathrm{C}_4^13^1x=12x$.
	}
\end{ex}
\begin{ex}%[0D8H3-4]
	Tìm hệ số của $x^2y^2$ trong khai triển nhị thức Newtơn của $(x+2y)^4$.
	\choice
	{$32$}
	{$8$}
	{\True $24$}
	{$16$}
	\loigiai{
		Ta có $(x+2y)^4=\sum\limits_{k=0}^4\mathrm{C}_4^kx^{4-k}(2y)^k
		=\sum\limits_{k=0}^4\mathrm{C}_4^k\cdot 2^k\cdot x^{4-k}y^k$.\\
		Số hạng chứa $x^2y^2$ trong khai triển trên ứng với 
		$\heva{&4-k=2 \\&k=2}\Leftrightarrow k=2$.\\
		Vậy hệ số của $x^2y^2$ trong khai triển của $(x+2y)^4$ là $\mathrm{C}_4^2\cdot 2^2=24$.
	}
\end{ex}
\begin{ex}%[0D8V3-4]
	Tìm số hạng chứa $x^2$ trong khai triển nhị thức $P(x)=4x^2+x(x-2)^4$.
	\choice
	{$28x^2$}
	{\True $-28x^2$}
	{$-24x^2$}
	{$24x^2$}
	\loigiai{
		Ta có $P(x)=4x^2+x(x-2)^4=4x^2+x\sum\limits_{k=0}^4\mathrm{C}_4^kx^{4-k}(-2)^k
		=4x^2+\sum\limits_{k=0}^4\mathrm{C}_4^k(-2)^kx^{5-k}$.\\
		Số hạng chứa $x^2$ (ứng với $k=3$) trong khai triển $P(x)$ là $\left[4+\mathrm{C}_4^3(-2)^3\right]x^2=-28x^2$.
	}
\end{ex}
\begin{ex}%[0D8V3-4]
	Gọi $n$ là số nguyên dương thỏa mãn $\mathrm{A}_n^3+2\mathrm{A}_n^2=48$. Tìm hệ số của $x^3$ trong khai triển nhị thức Newtơn của $(1-3x)^n$.
	\choice
	{\True $-108$}
	{$81$}
	{$54$}
	{$-12$}
	\loigiai{
		Điều kiện $n\ge 3$; $n\in \mathbb{N}$. Khi đó, ta có
		\begin{eqnarray*}
			\mathrm{A}_n^3+2\mathrm{A}_n^2=48 &\Leftrightarrow& \dfrac{n!}{(n-3)!}+2\cdot \dfrac{n!}{(n-2)!}=48 \\
			&\Leftrightarrow& n(n-1)(n-2)+2\cdot n(n-1)=48 \\
			&\Leftrightarrow& n^3-n^2-48=0 \Leftrightarrow n=4 (\text{nhận}).
		\end{eqnarray*}
		Ta có $(1-3x)^4=\sum\limits_{k=0}^4\mathrm{C}_4^k\left(-3x\right)^k
		=\sum\limits_{k=0}^4\mathrm{C}_4^k\left(-3\right)^kx^k$.\\
		Hệ số của $x^3$ trong khai triển trên ứng với $k=3$.\\
		Vậy hệ số của $x^3$ trong khai triển $(1-3x)^4$ là $\mathrm{C}_4^3\cdot \left(-3\right)^3=-108$.
	}
\end{ex}
\begin{ex}%[0D8V3-4]
	Tìm số hạng không chứa $x$ trong khai triển nhị thức Newtơn của $\left(\dfrac{1}{x}+x^3\right)^4$.
	\choice
	{$1$}
	{\True $4$}
	{$6$}
	{$12$}
	\loigiai{
	Ta có $\left(\dfrac{1}{x}+x^3\right)^4=\sum\limits_{k=0}^4\mathrm{C}_4^k\left(\dfrac{1}{x}\right)^{4-k}\left(x^3\right)^k=\sum\limits_{k=0}^4\mathrm{C}_4^kx^{4k-4}$.\\
		Số hạng không chứa $x$ trong khai triển trên ứng với $x^0$ hay $4k-4=0\Leftrightarrow k=1$.\\
		Vậy số hạng không chứa $x$ trong khai triển $\left(\dfrac{1}{x}+x^3\right)^4$ là $\mathrm{C}_4^1=4$.
	}
\end{ex}
\begin{ex}%[0D8H3-2]
	Viết khai triển theo công thức nhị thức Newtơn $(x-1)^5$.
	\choice
	{$x^5+5x^4+10x^3+10x^2+5x+1$}
	{$x^5-5x^4-10x^3+10x^2-5x+1$}
	{\True $x^5-5x^4+10x^3-10x^2+5x-1$}
	{$5x^5+10x^4+10x^3+5x^2+5x+1$}
	\loigiai{
		$\left(x-1\right)^5=\mathrm{C}_5^0x^5+\mathrm{C}_5^1x^4+\mathrm{C}_5^2x^3+\mathrm{C}_5^3x^2+\mathrm{C}_5^4x+\mathrm{C}_5^5=x^5+5x^4+10x^3+10x^2+5x+1$.
	}
\end{ex}
\begin{ex}%[0D8H3-2]
	Viết khai triển theo công thức nhị thức Newtơn $(x-y)^5$.
	\choice
	{\True $x^5-5x^4y+10x^3y^2-10x^2y^3+5xy^4-y^5$}
	{$x^5+5x^4y+10x^3y^2+10x^2y^3+5xy^4+y^5$}
	{$x^5-5x^4y-10x^3y^2-10x^2y^3-5xy^4+y^5$}
	{$x^5+5x^4y-10x^3y^2+10x^2y^3-5xy^4+y^5$}
	\loigiai{
		$\begin{aligned}
			\text{Ta có }
			(x-y)^5&=\mathrm{C}_5^0x^5+\mathrm{C}_5^1x^4(-y)+\mathrm{C}_5^2x^3(-y)^2+\mathrm{C}_5^3x^2(-y)^3+\mathrm{C}_5^4x(-y)^4+\mathrm{C}_5^5(-y)^5 \\
			&=x^5-5x^4y+10x^3y^2-10x^2y^3+5xy^4-y^5.
		\end{aligned}$
	}
\end{ex}
\begin{ex}%[0D8H3-2]
	Khai triển của nhị thức $(x-2)^5$.
	\choice
	{$x^5-100x^4+400x^3-800x^2+800x-32$}
	{$5x^5-10x^4+40x^3-80x^2+80x-32$}
	{\True $x^5-10x^4+40x^3-80x^2+80x-32$}
	{$x^5+10x^4+40x^3+80x^2+80x+32$}
	\loigiai{
		$\begin{aligned}
			\text{Ta có }
			(x-2)^5&=\mathrm{C}_5^0x^5+\mathrm{C}_5^1x^4(-2)+\mathrm{C}_5^2x^3(-2)^2+\mathrm{C}_5^3x^2(-2)^3+\mathrm{C}_5^4x(-2)^4+\mathrm{C}_5^5(-2)^5 \\
			&=x^5-10x^4+40x^3-80x^2+80x-32.
		\end{aligned}$
	}
\end{ex}
\begin{ex}%[0D8H3-2]
	Khai triển của nhị thức $(3x+4)^5$ là
	\choice
	{$x^5+1620x^4+4320x^3+5760x^2+3840x+1024$}
	{$243x^5+405x^4+4320x^3+5760x^2+3840x+1024$}
	{$243x^5-1620x^4+4320x^3-5760x^2+3840x-1024$}
	{\True $243x^5+1620x^4+4320x^3+5760x^2+3840x+1024$}
	\loigiai{
		$\begin{aligned}
			\text{Ta có }
			(3x+4)^5&=\mathrm{C}_5^0(3x)^5+\mathrm{C}_5^1(3x)^4\cdot 4+\mathrm{C}_5^2(3x)^3\cdot 4^2+\mathrm{C}_5^3(3x)^2\cdot 4^3 +\mathrm{C}_5^4(3x)^1\cdot 4^4 +\mathrm{C}_5^5\cdot 4^5 \\
			&=243x^5+1620x^4+4320x^3+5760x^2+3840x+1024.
		\end{aligned}$
	}
\end{ex}
\begin{ex}%[0D8H3-2]
	Khai triển của nhị thức $(1-2x)^5$ là
	\choice
	{$5-10x+40x^2-80x^3-80x^4-32x^5$}
	{$1+10x+40x^2-80x^3-80x^4-32x^5$}
	{\True $1-10x+40x^2-80x^3+80x^4-32x^5$}
	{$1+10x+40x^2+80x^3+80x^4+32x^5$}
	\loigiai{
		$\begin{aligned}
			\text{Ta có }
			(1-2x)^5&=\mathrm{C}_5^0+\mathrm{C}_5^1(-2x)^1+\mathrm{C}_5^2(-2x)^2+\mathrm{C}_5^3(-2x)^3+\mathrm{C}_5^4(-2x)^4+\mathrm{C}_5^5(-2x)^5 \\
			&=1-10x+40x^2-80x^3+80x^4-32x^5.
		\end{aligned}$
	}
\end{ex}
\Closesolutionfile{ans}
\indapan{6}{ans}

\cauds
\Opensolutionfile{ans}[ans/0-C8B3-1]
\begin{ex}%[0D8H3-4]
	Khai triển $\left(x+\sqrt{2}\right)^4$. Các mệnh đề sau đúng hay sai?
	\choiceTF
	{\True Hệ số của $x^2$ là $12$}		
	{Khai triển trên có $4$ số hạng}		
	{\True Hệ số lớn nhất của khai triển là $8\sqrt{2}$}		
	{\True Số hạng không chứa $x$ trong khai triển trên bằng $4$}
\loigiai{
	Ta có \\
	$\begin{aligned}
	\left(x+\sqrt{2}\right)^4&=\mathrm{C}_4^0 x^4+\mathrm{C}_4^1 x^3\left(\sqrt{2}\right)+\mathrm{C}_4^2 x^2\left(\sqrt{2}\right)^2+\mathrm{C}_4^3 x\left(\sqrt{2}\right)^3+\mathrm{C}_4^4\left(\sqrt{2}\right)^4 \\
	&=x^4+4 \sqrt{2} x^3+12 x^2+8 \sqrt{2} x+4.
	\end{aligned}$
}
\end{ex}
\begin{ex}%[0D8V3-4]
	Khai triển $(x+2y)^3+(2x-y)^3$. Khi đó
	\choiceTF
	{\True Hệ số của của $x^3$ là $9$}
	{\True Hệ số của của $y^3$ là $7$}
	{Hệ số của $x^2y$ là $6$}
	{Tổng các hệ số của số hạng mà lũy thừa của $x$ lớn hơn lũy thừa của $y$ bằng $-3$}
\loigiai{
	Ta có\\
	$\begin{aligned}
		(x+2y)^3+(2x-y)^3&=\mathrm{C}_3^0x^3+\mathrm{C}_3^1x^2(2y)+\mathrm{C}_3^2x(2y)^2+\mathrm{C}_3^3(2y)^3+\mathrm{C}_3^0(2x)^3+\mathrm{C}_3^1(2x)^2(-y) \\
		&\hspace*{7.5cm} +\mathrm{C}_3^2(2x)(-y)^2+\mathrm{C}_3^3(-y)^3 \\
		&=x^3+6x^2y+12xy^2+8y^3+8x^3-12x^2y+6xy^2-y^3 \\
			&=9x^3-6x^2y+18xy^2+7y^3.
	\end{aligned}$\\
	Có hai số hạng mà lũy thừa của $x$ lớn hơn lũy thừa của $y$ là $9x^3$ và $-6x^2y$. \\
	Tổng hệ số của chúng là $9+(-6)=3$.
}	
\end{ex}
\begin{ex}%[0D8H3-4]
	Khai triển $\left(x+\dfrac{1}{x}\right)^4$. Khi đó
	\choiceTF
	{Hệ số của $x^2$ là $\dfrac{1}{4}$} 
	{\True Số hạng không chứa $x$ là $6$} 
	{\True Hệ số của $x^4$ là $1$}
	{\True Sau khi khai triển, biểu thức có $5$ số hạng}
\loigiai{
	$\begin{aligned}
		\text{Ta có }\left(x+\dfrac{1}{x}\right)^4
		&=\mathrm{C}_4^0 x^4+\mathrm{C}_4^1 x^3\left(\dfrac{1}{x}\right)+\mathrm{C}_4^2 x^2\left(\dfrac{1}{x}\right)^2+\mathrm{C}_4^3 x\left(\dfrac{1}{x}\right)^3+\mathrm{C}_4^4\left(\dfrac{1}{x}\right)^4 \\
		&=x^4+4 x^2+6+\dfrac{4}{x^2}+\dfrac{1}x^4.
	\end{aligned}$
}
\end{ex}
\begin{ex}%[0D8V3-5]
	Khai triển $(x+1)^5$. Khi đó
	\choiceTF
	{\True Hệ số của $x^4$ là $5$}
	{\True Số hạng không chứa $x$ là $1$}
	{$\mathrm{C}_5^0+\mathrm{C}_5^1+\mathrm{C}_5^2+\mathrm{C}_5^3+\mathrm{C}_5^4+\mathrm{C}_5^5=3^5$}
	{\True $32\mathrm{C}_5^0+16\mathrm{C}_5^1+8\mathrm{C}_5^2+4\mathrm{C}_5^3+2\mathrm{C}_5^4+\mathrm{C}_5^5=3^5$}
\loigiai{
	$\begin{aligned}
		\text{Ta có }(x+1)^5&=\mathrm{C}_5^0x^5+\mathrm{C}_5^1x^4+\mathrm{C}_5^2x^3+\mathrm{C}_5^3x^2+\mathrm{C}_5^4x+\mathrm{C}_5^5 \hfill \left( * \right)\\ 
		&=1+5x+10x^2+10x^3+5x^4+x^5.
	\end{aligned}$\\
	Từ khai triển $(*)$ trong câu $a)$, thay $x=1$, ta được: 
	$$\begin{aligned}
		(1+1)^5&=\mathrm{C}_5^0 \cdot 1^5+\mathrm{C}_5^1 \cdot 1^4+\mathrm{C}_5^2 \cdot 1^3+\mathrm{C}_5^3 \cdot 1^2+\mathrm{C}_5^4 \cdot 1+\mathrm{C}_5^5 \\ &=\mathrm{C}_5^0+\mathrm{C}_5^1+\mathrm{C}_5^2+\mathrm{C}_5^3+\mathrm{C}_5^4+\mathrm{C}_5^5.
	\end{aligned}$$
	Vậy $\mathrm{C}_5^0+\mathrm{C}_5^1+\mathrm{C}_5^2+\mathrm{C}_5^3+\mathrm{C}_5^4+\mathrm{C}_5^5=2^5$.\\
	Từ khai triển $(*)$ của câu $a)$, thay $x=2$, ta được:
	$$\begin{aligned}
		(2+1)^5&=\mathrm{C}_5^0\cdot 2^5+\mathrm{C}_5^1\cdot 2^4+\mathrm{C}_5^2\cdot 2^3+\mathrm{C}_5^3\cdot 2^2+\mathrm{C}_5^4\cdot 2+\mathrm{C}_5^5 \\ 
		&=32\mathrm{C}_5^0+16\mathrm{C}_5^1+8\mathrm{C}_5^2+4\mathrm{C}_5^3+2\mathrm{C}_5^4+\mathrm{C}_5^5=S. 
	\end{aligned}$$
	Vậy $S=3^5$.
}
\end{ex}
\Closesolutionfile{ans}
\indapan{3}{ans/0-C8B3-1}

\Opensolutionfile{ans}[ans/ans-0-B1-KQ]
\caukq
\begin{ex}%[0D8V3-4]
	Tìm số hạng chứa $x^3$ trong khai triển của đa thức $x(2x+1)^4+(x+2)^5$.
	\shortans{$64x^3$}
\loigiai{
	\begin{itemize}
		\item Ta có $x(2x+1)^4=x\sum\limits_{k=0}^4 C_4^k2^{4-k}x^{4-k}$. \\
		Số hạng chứa $x^3$ tương ứng với $4-k=2\Rightarrow k=2$.
		\item Ta có $(x+2)^5=\sum\limits_{k=0}^5 C_5^k\cdot 2^kx^{5-k}$. \\
		Số hạng chứa $x^3$ tương ứng với $5-k=3\Rightarrow k=2$.
	\end{itemize}
	Vậy số hạng cần tìm là $C_4^2\cdot 2^2x^3+C_5^2\cdot 2^2x^3=64x^3$.
}
\end{ex}
\begin{ex}%[0D8V3-5]
	Lớp $10A$ đề nghị các tổ chọn thành viên để tập kịch. Tổ I phải chọn ít nhất một thành viên để tham gia đội kịch của lớp. Hỏi tổ I có bao nhiêu cách chọn thành viên để tập kịch? Biết rằng tổ I có $5$ người.
	\shortans{$31$}
\loigiai{
	Vì tổ I phải chọn ít nhất một thành viên để tham gia đội kịch nên số cách chọn thành viên của tổ I là
	\begin{eqnarray*}
		&&\mathrm{C}_5^1+\mathrm{C}_5^2+\mathrm{C}_5^3+\mathrm{C}_5^4+\mathrm{C}_5^5=(1+1)^5-\mathrm{C}_5^0=2^5-1=31.
	\end{eqnarray*}
} 
\end{ex}
\begin{ex}%[0D8V3-5]
	Cho tập hợp $A=\{1; 2; 3; 4; 5; 6\}$. Hỏi tập $A$ có bao nhiêu tập hợp con?
	\shortans{$2^6$}
\loigiai{
	Số tập con có $k$ phần tử của $A$ là $\mathrm{C}_6^k$ với $k=\overline{0,6}$.\\
	Vậy tổng số tập con của $A$ là 
	$\sum\limits_{k=0}^6\, \mathrm{C}_6^k = {\mathrm{C}_6^0+\mathrm{C}_6^1+\mathrm{C}_6^2+\mathrm{C}_6^3+\mathrm{C}_6^4+\mathrm{C}_6^5+\mathrm{C}_6^6=T}$. \\
	Theo khai triển nhị thức Newtơn, ta có:
	$$\begin{aligned}
		(1+x)^6=\mathrm{C}_6^0+\mathrm{C}_6^1x+\mathrm{C}_6^2x^2+\mathrm{C}_6^3x^3+\mathrm{C}_6^4x^4+\mathrm{C}_6^5x^5+\mathrm{C}_6^6x^6.
	\end{aligned}$$
	Thay $x=1$, ta được $(1+1)^6=\mathrm{C}_6^0+\mathrm{C}_6^1+\mathrm{C}_6^2+\mathrm{C}_6^3+\mathrm{C}_6^4+\mathrm{C}_6^5+\mathrm{C}_6^6$ hay $T=2^6$.\\
	Vậy số tập con của tập $A$ là $2^6$.
}
\end{ex}
\begin{ex}%[0D8V3-4]
	Cho $n$ là số nguyên dương thỏa mãn ${\mathrm{C}_n^1+\mathrm{C}_n^2=15}$. Tìm số hạng không chứa $x$ trong khai triển: $\left(x+\dfrac{2}{x^4}\right)^n$.
	\shortans{$10$}
\loigiai{
	Điều kiện $n \geq 2, n \in \mathbb{N}^*$. Ta có
	\begin{eqnarray*}
		\mathrm{C}_n^1+\mathrm{C}_n^2=15 
		&\Leftrightarrow& n+\dfrac{n(n-1)}{2}=15 \\
		&\Leftrightarrow& n^2+n-30=0 \\
		&\Leftrightarrow& \hoac{&n=5 \\ &n=-6} \Rightarrow n=5.
	\end{eqnarray*} 
	Khi đó $\left(x+\dfrac{2}{x^4}\right)^5=\sum\limits_{k=0}^5 \mathrm{C}_5^k \cdot 2^k x^{5-k} \cdot \left(\dfrac{1}{x^4}\right)^k =\sum\limits_{k=0}^5 \mathrm{C}_5^k \cdot 2^kx^{5-5k}$.\\
	Số hạng không chứa $x$ tương ứng $5-5k=0 \Leftrightarrow k=1$. \\
	Suy ra số hạng không chứa $x$ là $\mathrm{C}_5^1 \cdot 2^1=10$.
}
\end{ex}
\begin{ex}%[0D8V3-6]
	Cho khai triển $(1+2x)^n=a_0+a_1x+a_2x^2+\ldots+a_nx^n$ thỏa mãn $a_0+8 a_1=2 a_2+1$. Tìm giá trị của số nguyên dương $n$.
	\shortans{$5$}
\loigiai{
	Ta có $(1+2x)^n=\sum\limits_{k=0}^n2^k\mathrm{C}_n^kx^k$; $(k \in \mathbb{N})$. Suy ra $a_k=2^k \mathrm{C}_n^k$. \\
	Thay $a_0=2^0\mathrm{C}_n^0=1$, $a_1=2 \mathrm{C}_n^1$, $a_2=4 \mathrm{C}_n^2$ vào giả thiết ta có
	\begin{eqnarray*}
		1+16\mathrm{C}_n^1=8\mathrm{C}_n^2+1 
		&\Leftrightarrow& 2\mathrm{C}_n^1=\mathrm{C}_n^2
		\Leftrightarrow 2\dfrac{n!}{(n-1)!}=\dfrac{n!}{(n-2)!2!} \\
		&\Leftrightarrow& 2n=\dfrac{n(n-1)}{2}
		\Leftrightarrow n^2-5 n=0 \Leftrightarrow \hoac{&n=0 \\ &n=5.}
	\end{eqnarray*}
	Do $n$ là số nguyên dương nên $n=5$.
}
\end{ex}
\begin{ex}%[0D8V3-4]
	Tìm hệ số của $x^{10}$ trong khải triển thành đa thức của $\left(1+x+x^2+x^3\right)^5$.
	\shortans{$101$}
\loigiai{
	Ta có $\left(1+x+x^2+x^3\right)^5
	=\left[(1+x) \cdot\left(1+x^2\right)\right]^5
	=(1+x)^5\cdot \left(1+x^2\right)^5$.\\
	Xét khai triển $(1+x)^5\cdot \left(1+x^2\right)^5
	=\sum\limits_{k=0}^5 \mathrm{C}_5^k x^k\cdot \sum_{l=0}^5 \mathrm{C}_5^l x^{2l}
	=\sum\limits_{k=0}^5\left(\mathrm{C}_5^k\cdot \sum_{l=0}^5 \mathrm{C}_5^l \cdot x^{k+2l}\right)$. \\
	Số hạng chứa $x^{10}$ tương ứng với $k, l$ thỏa mãn $k+2l=10 \Leftrightarrow k=10-2l$. \\
	Kết hợp với điều kiện, ta có hệ:
	$\heva{&k=10-2l \\ &0 \leq k \leq 5\, (k \in \mathbb{N}) \\ &0 \leq l \leq 5\, (l \in \mathbb{N})} \Leftrightarrow (k, l) \in \{(0; 5),(2; 4),(4; 3)\}$.\\
	Vậy hệ số của $x^{10}$ bằng tổng các $\mathrm{C}_5^k \cdot \mathrm{C}_5^l$ thỏa mãn
	$\mathrm{C}_5^0\cdot \mathrm{C}_5^5+\mathrm{C}_5^2\cdot \mathrm{C}_5^4+\mathrm{C}_5^4\cdot \mathrm{C}_5^3=101$.
}
\end{ex}
\Closesolutionfile{ans}
\indapan{6}{ans/ans-0-B1-KQ}
\newpage
\subsection{Đề 2}
\Opensolutionfile{ans}[ans]
\caulc
\begin{ex}%[0D8H3-2]
	Đa thức $P(x)=32x^5-80x^4+80x^3-40x^2+10x-1$ là khai triển của nhị thức nào dưới đây?
	\choice
	{$(1-2x)^5$}
	{$(1+2x)^5$}
	{\True $(2x-1)^5$}
	{$(x-1)^5$}
	\loigiai{
		Ta có $P(x)=32x^5-80x^4+80x^3-40x^2+10x-1=(2x-1)^5$.
	}
\end{ex}
\begin{ex}%[0D8H3-2]
	Khai triển nhị thức $(2x+y)^5$. Ta được kết quả là
	\choice
	{$32x^5+16x^4y+8x^3y^2+4x^2y^3+2xy^4+y^5$}
	{\True $32x^5+80x^4y+80x^3y^2+40x^2y^3+10xy^4+y^5$}
	{$2x^5+10x^4y+20x^3y^2+20x^2y^3+10xy^4+y^5$}
	{$32x^5+10000x^4y+80000x^3y^2+400x^2y^3+10xy^4+y^5$}
	\loigiai{
		$\begin{aligned}
			\text{Ta có }
			(2x+y)^5&=\mathrm{C}_5^0(2y)^5+\mathrm{C}_5^1(2y)^4y+\mathrm{C}_5^2(2y)^3y^2+\mathrm{C}_5^3(2y)^2y^3+\mathrm{C}_5^4(2y)y^4+\mathrm{C}_5^5y^5 \\
			&=32x^5+80x^4y+80x^3y^2+40x^2y^3+10xy^4+y^5.
		\end{aligned}$
	}
\end{ex}
\begin{ex}%[0D8H3-2]
	Đa thức $P(x)=x^5-5x^4y+10x^3y^2-10x^2y^3+5xy^4-y^5$ là khai triển của nhị thức nào dưới đây?
	\choice
	{\True $(x-y)^5$}
	{$(x+y)^5$}
	{$(2x-y)^5$}
	{$(x-2y)^5$}
	\loigiai{
		Ta có $P(x)=x^5-5x^4y+10x^3y^2-10x^2y^3+5xy^4-y^5=(x-y)^5$.
	}
\end{ex}
\begin{ex}%[0D8V3-2]
	Khai triển của nhị thức $\left(x-\dfrac{1}{x}\right)^5$ là
	\choice
	{$x^5+5x^3+10x+\dfrac{10}{x}+\dfrac{5}{x^3}+\dfrac{1}{x^5}$}
	{\True $x^5-5x^3+10x-\dfrac{10}{x}+\dfrac{5}{x^3}-\dfrac{1}{x^5}$}
	{$5x^5-10x^3+10x-\dfrac{10}{x}+\dfrac{5}{x^3}-\dfrac{1}{x^5}$}
	{$5x^5+10x^3+10x+\dfrac{10}{x}+\dfrac{5}{x^3}+\dfrac{1}{x^5}$}
	\loigiai{
		$\begin{aligned}
			\text{Ta có }
			\left(x-\dfrac{1}{x}\right)^5&=\mathrm{C}_5^0\cdot x^5+\mathrm{C}_5^1\cdot x^4\cdot \left(\dfrac{-1}{x}\right)^1+\mathrm{C}_5^2x^3\left(\dfrac{-1}{x}\right)^2+\mathrm{C}_5^3x^2\left(\dfrac{-1}{x}\right)^3\\
			&\hspace*{8cm}+\mathrm{C}_5^4x^1\left(\dfrac{-1}{x}\right)^4+\mathrm{C}_5^5\left(\dfrac{-1}{x}\right)^5 \\
			&=x^5-5x^3+10x-\dfrac{10}{x}+\dfrac{5}{x^3}-\dfrac{1}{x^5}\cdot
		\end{aligned}$
	}
\end{ex}
\begin{ex}%[0D8H3-2]
	Khai triển của nhị thức $(xy+2)^5$ là
	\choice
	{\True $x^5y^5+10x^4y^4+40x^3y^3+80x^2y^2+80xy+32$}
	{$5x^5y^5+10x^4y^4+40x^3y^3+80x^2y^2+80xy+32$}
	{$x^5y^5+100x^4y^4+400x^3y^3+80x^2y^2+80xy+32$}
	{$x^5y^5-10x^4y^4+40x^3y^3-80x^2y^2+80xy-32$}
	\loigiai{
		$\begin{aligned}
			\text{Ta có }
			(xy+2)^5&=\mathrm{C}_5^0(xy)^5+\mathrm{C}_5^1(xy)^4\cdot 2^1+\mathrm{C}_5^2(xy)^3\cdot 2^2+\mathrm{C}_5^3(xy)^2\cdot 2^3+\mathrm{C}_5^4(xy)^1\cdot 2^4+\mathrm{C}_5^5\cdot 2^5 \\
			&=x^5y^5+10x^4y^4+40x^3y^3+80x^2y^2+80xy+32.
		\end{aligned}$
	}
\end{ex}
\begin{ex}%[0D8H3-4]
	Trong khai triển $\left(2a-b\right)^5$, hệ số của số hạng thứ $3$ bằng
	\choice
	{$-80$}
	{\True $80$}
	{$-10$}
	{$10$}
	\loigiai{
		$\begin{aligned}
			\text{Ta có }
			\left(2a-b\right)^5&=(2a)^5-5(2a)^4b+10(2a)^3b^2-10(2a)^2b^3+5(2a)b^4-b^5 \\
			&=32a^5-80a^4b+80a^3b^2-40a^2b^3+10ab^4-b^5.
		\end{aligned}$
	}
\end{ex}
\begin{ex}%[0D8H3-4]
	Tìm hệ số của đơn thức $a^3b^2$ trong khai triển nhị thức $(a+2b)^5$.
	\choice
	{$160$}
	{$80$}
	{$20$}
	{\True $40$}
	\loigiai{
		$\begin{aligned}
			\text{Ta có }
			(a+2b)^5&=a^5+5a^4(2b)+10a^3(2b)^2+10a^2(2b)^3+5a(2b)^4+(2b)^5 \\
			& =a^5+10a^4b+40a^3b^2+80a^2b^3+80ab^4+32b^5.
		\end{aligned}$\\
		Suy ra hệ số của $a^3b^2$ trong khai triển trên là $40$.
	}
\end{ex}
\begin{ex}%[0D8H3-4]
	Số hạng chính giữa trong khai triển $\left(3x+2y\right)^4$ là
	\choice
	{$\mathrm{C}_4^2x^2y^2$}
	{$6(3x)^2(2y)^2$}
	{$6\mathrm{C}_4^2x^2y^2$}
	{\True $36\mathrm{C}_4^2x^2y^2$}
	\loigiai{
		Ta có $\left(3x+2y\right)^4=(3x)^4+4(3x)^3(2y)+6(3x)^2(2y)^2+4(3x)(2y)^3+(2y)^4$. \\
		Suy ra hệ số chính giữa trong khai triển trên là $6(3x)^2(2y)^2=36\mathrm{C}_4^2x^2y^2$.
	}
\end{ex}
\begin{ex}%[0D8V3-4]
	Biết $\left(1+\sqrt[3]{2}\right)^4=a_0+a_1\sqrt[3]{2}+a_2\sqrt[3]{4}$. Tính $a_1a_2$.
	\choice
	{$a_1a_2=24$}
	{$a_1a_2=8$}
	{$a_1a_2=54$}
	{\True $a_1a_2=36$}
	\loigiai{
		$\begin{aligned}
			\text{Ta có }
			\left(1+\sqrt[3]{2}\right)^4&=1^4+4\cdot 1^3\left(\sqrt[3]{2}\right)^1+6\cdot 1^2\left(\sqrt[3]{2}\right)^2+4\cdot 1^1\left(\sqrt[3]{2}\right)^3+\left(\sqrt[3]{2}\right)^4 \\
			&=1+4\sqrt[3]{2}+6\sqrt[3]{4}+8+2\sqrt[3]{2}
			=9+6\sqrt[3]{2}+6\sqrt[3]{4}.
		\end{aligned}$\\
		Suy ra $a_1a_2=6\cdot 6=36$.
	}
\end{ex}
\begin{ex}%[0D8H3-4]
	Số hạng chứa $\sqrt{x}$ trong khai triển $\left(\sqrt{x}+\dfrac{2}{x}\right)^4$, $x>0$ là số hạng thứ mấy?
	\choice
	{$5$}
	{$3$}
	{\True $2$}
	{$4$}
	\loigiai{
		$\begin{aligned}
			\text{Ta có }
			\left(\sqrt{x}+\dfrac{2}{x}\right)^4&=\left(\sqrt{x}\right)^4+4\left(\sqrt{x}\right)^3\left(\dfrac{2}{x}\right)+6\left(\sqrt{x}\right)^2\left(\dfrac{2}{x}\right)^2+4\left(\sqrt{x}\right)\left(\dfrac{2}{x}\right)^3+\left(\dfrac{2}{x}\right)^4 \\
			&=x^2+8\sqrt{x}+24\dfrac{1}{x}+32\dfrac{\sqrt{x}}{x^3}+16\dfrac{1}{x^4}\cdot
		\end{aligned}$\\
		Số hạng chứa $\sqrt{x}$ trong khai triển trên ứng với số hạng thứ $2$.
	}
\end{ex}
\begin{ex}%[0D8H3-4]
	Tìm số hạng không chứa $x$ trong khai triển của nhị thức $\left(x^3-\dfrac{1}{x^2}\right)^5$.
	\choice
	{\True $-10$}
	{$-5$}
	{$10$}
	{$5$}
	\loigiai{
		$\begin{aligned}
			\text{Ta có }
			\left(x^3-\dfrac{1}{x^2}\right)^5&=\left(x^3\right)^5-5\left(x^3\right)^4\left(\dfrac{1}{x^2}\right)+10\left(x^3\right)^3\left(\dfrac{1}{x^2}\right)^2-10\left(x^3\right)^2\left(\dfrac{1}{x^2}\right)^3 \\
			&\hspace*{8cm}+5\left(x^3\right)\left(\dfrac{1}{x^2}\right)^4-\left(\dfrac{1}{x^2}\right)^5 \\
			&=x^{15}-5x^{10}+10x^5-10+\dfrac{5}{x^5}-\dfrac{1}{x^{10}}\cdot
		\end{aligned}$\\
		Số hạng không chứa $x$ trong khai triển là $-10$.
	}
\end{ex}
\begin{ex}%[0D8V3-5]
	Rút gọn
	$M=\mathrm{C}_4^0a^4+\mathrm{C}_4^1a^3(1-a)+\mathrm{C}_4^2a^2(1-a)^2+\mathrm{C}_4^3a(1-a)^3+\mathrm{C}_4^4(1-a)^4$.
	\choice
	{$M=a^4$}
	{$M=a$}
	{\True $M=1$}
	{$M=-1$}
	\loigiai{
		$\begin{aligned}
			\text{Ta có }
			M&=\mathrm{C}_4^0a^4+\mathrm{C}_4^1a^3(1-a)+\mathrm{C}_4^2a^2(1-a)^2+\mathrm{C}_4^3a(1-a)^3+\mathrm{C}_4^4(1-a)^4 \\
			&=\left[a+(1-a)\right]^4=1.
		\end{aligned}$
	}
\end{ex}
\begin{ex}%[0D8H3-4]
	Cho $\left(1-2x\right)^n=a_0+a_1x+a_2x^2+\ldots+a_nx^n$. Tìm $a_4$ biết $a_0+a_1+a_2=31$.
	\choice
	{\True $80$}
	{$-80$}
	{$40$}
	{$-40$}
	\loigiai{
		$\begin{aligned}
			\text{Ta có }
			(1-2x)^n&=\mathrm{C}_n^01^n(-2x)^0+\mathrm{C}_n^11^{n-1}(-2x)+\mathrm{C}_n^21^{n-2}(-2x)^2+\ldots \\
			&=1-2\mathrm{C}_n^1x+4\mathrm{C}_n^2x^2+\ldots
		\end{aligned}$
		\\
		Vậy $a_0=1$; $a_1=-2\mathrm{C}_n^1$; $a_2=4\mathrm{C}_n^2$.\\
		Theo bài ra $a_0+a_1+a_2=31$ nên ta có:
		\begin{eqnarray*}
			1-2\mathrm{C}_n^1+4\mathrm{C}_n^2=31
			&\Leftrightarrow& 1-2\dfrac{n!}{1!\left(n-1\right)!}+4\dfrac{n!}{2!\left(n-2\right)!}=31 \\
			&\Leftrightarrow& 1-2n+2n\left(n-1\right)=31 \\
			&\Leftrightarrow& 2n^2-4n-30=0 \\
			&\Leftrightarrow& n^2-2n-15=0 \Rightarrow n=5.
		\end{eqnarray*}
		Từ đó ta có $a_4=\mathrm{C}_5^4(-2)^4=80$.
	}
\end{ex}
\begin{ex}%[0D8V3-4]
	Biết hệ số của $x^2$ trong khai triển của $(1-3x)^n$ là $90$. Khi đó ta có $3n^4$ bằng
	\choice
	{$7203$}
	{\True $1875$}
	{$1296$}
	{$6561$}
	\loigiai{
		Nhị thức $(1-3x)^n$ có $T_{k+1}=\mathrm{C}_n^k\left(-3x\right)^k=\left(-3\right)^k\mathrm{C}_n^kx^k$.\\
		$\Rightarrow $ hệ số của $x^2$ trong khai triển của $(1-3x)^n$ ứng với $k=2$.\\
		Khi đó $\left(-3\right)^2\mathrm{C}_n^2=90\Leftrightarrow 9\dfrac{n\left(n-1\right)}{2}=90\Leftrightarrow n\left(n-1\right)=20\Leftrightarrow \hoac{&n=-4 \\&n=5}$ $\Rightarrow 3n^4=1875$.
	}
\end{ex}
\begin{ex}%[0D8V3-4]
	Tìm hệ số của $x^2$ trong khai triển $f(x)=\left(x^3+\dfrac{1}{x^2}\right)^n$, với $x>0$, \break biết: $\mathrm{C}_n^0+\mathrm{C}_n^1+\mathrm{C}_n^2=11$.
	\choice
	{$20$}
	{\True $6$}
	{$7$}
	{$15$}
	\loigiai{
		Ta có $\mathrm{C}_n^0+\mathrm{C}_n^1+\mathrm{C}_n^2=11\Leftrightarrow 1+n+\dfrac{n\left(n-1\right)}{2}=11$ $\Leftrightarrow \hoac{&n=4 \\&n=-5}$.\\
		Nhị thức $f(x)=\left(x^3+\dfrac{1}{x^2}\right)^4$ có $T_{k+1}=\mathrm{C}_4^k\left(x^3\right)^{4-k}\left(\dfrac{1}{x^2}\right)^k=\mathrm{C}_4^kx^{12-5k}$.\\
		Số hạng chứa $x^2$ trong khai triển ứng với số mũ của $x$ là $12-5k=2\Leftrightarrow k=2$.\\
		Vậy hệ số của $x^2$ trong khai triển là $\mathrm{C}_4^2=6$.
	}
\end{ex}
\begin{ex}%[0D8V3-4]
	Tìm hệ số của $x^2$ trong khai triển: $f(x)=\left(x^3+\dfrac{2}{x^2}\right)^n$, với $x>0$, biết tổng ba hệ số đầu của $x$ trong khai triển bằng $33$.
	\choice
	{$34$}
	{\True $24$}
	{$6$}
	{$12$}
	\loigiai{
		Ta có $\mathrm{C}_n^0+2\mathrm{C}_n^1+4\mathrm{C}_n^2=33\Rightarrow n=4$. \\
		Nhị thức $f(x)=\left(x^3+\dfrac{2}{x^2}\right)^4$ có $T_{k+1}=\mathrm{C}_4^k\left(x^3\right)^{4-k}\left(\dfrac{2}{x^2}\right)^k=2^k\mathrm{C}_4^kx^{12-5k}$.\\
		Số hạng chứa $x^2$ trong khai triển ứng với số mũ của $x$ là $12-5k=2\Leftrightarrow k=2$.\\
		Vậy hệ số của $x^2$ trong khai triển là $2^2\mathrm{C}_4^2=24$.
	}
\end{ex}
\Closesolutionfile{ans}
\indapan{8}{ans}

\cauds
\Opensolutionfile{ans}[ans/0C8B3-1]
\begin{ex}%[0D8H3-4]
	Khai triển $P=\left(x-\sqrt{3}\right)^5$. Khi đó
	\choiceTF
	{Hệ số của $x^4$ trong khai triển là $5\sqrt{3}$}
	{\True Hệ số của $x^2$ trong khai triển là $-30\sqrt{3}$}
	{\True Hệ số của $x^3$ trong khai triển là $30$}
	{\True Hệ số của $x$ trong khai triển là $45$}
\loigiai{
	$\begin{aligned}
	\text{Ta có }
	P&=\left(x-\sqrt{3}\right)^5 \\
	&=\mathrm{C}_5^0 x^5+\mathrm{C}_5^1 x^4\left(-\sqrt{3}\right)+\mathrm{C}_5^2 x^3\left(-\sqrt{3}\right)^2+\mathrm{C}_5^3 x^2\left(-\sqrt{3}\right)^3 
	+\mathrm{C}_5^4x\left(-\sqrt{3}\right)^4+\mathrm{C}_5^5\left(-\sqrt{3}\right)^5 \\ 
	& =x^5-5\sqrt{3}x^4+30x^3-30\sqrt{3}x^2+45x-9\sqrt{3}.
	\end{aligned}$ \\
	Hệ số của $x^4$ trong khai triển là $-5\sqrt{3}$.
}
\end{ex}

\begin{ex}%[0D8H3-4]
	Khai triển $Q=(xy-1)^5$. Khi đó
	\choiceTF
	{\True Số hạng có chứa $x^2y^2$  là $-10x^2y^2$}
	{\True Hệ số của $x^4y^4$ trong khai triển là $-5$}
	{\True Hệ số của $x^3y^3$ trong khai triển là $10$}
	{Hệ số của $xy$ trong khai triển là $-10$}
\loigiai{
	$\begin{aligned}
		\text{Ta có }
		Q&=(xy-1)^5 \\
		&=\mathrm{C}_5^0(xy)^5+\mathrm{C}_5^1(xy)^4(-1)+\mathrm{C}_5^2(xy)^3(-1)^2
		+\mathrm{C}_5^3(xy)^2(-1)^3+\mathrm{C}_5^4(xy)(-1)^4+\mathrm{C}_5^5(-1)^5 \\ 
		& =x^5y^5-5x^4y^4+10x^3y^3-10x^2y^2+5xy-1.
	\end{aligned}$ \\
	Số hạng có chứa $x^2y^2$ trong khai triển là $-10x^2y^2$.
}
\end{ex}

\begin{ex}%[0D8H3-4]
	Khai triển $(1-x)^6$. Khi đó
	\choiceTF
	{\True Hệ số của $x^2$ trong khai triển là $\mathrm{C}_6^2$}
	{Hệ số của $x^3$ trong khai triển là $\mathrm{C}_6^3$}
	{\True Hệ số của $x^5$ trong khai triển là $-\mathrm{C}_6^5$}
	{$\mathrm{C}_6^0-\mathrm{C}_6^1+\mathrm{C}_6^2-\mathrm{C}_6^3+\mathrm{C}_6^4-\mathrm{C}_6^5+\mathrm{C}_6^6=1$}
\loigiai{
	$\begin{aligned}
		\text{Ta có }
		(1-x)^6=\mathrm{C}_6^0-\mathrm{C}_6^1x+\mathrm{C}_6^2x^2-\mathrm{C}_6^3x^3+\mathrm{C}_6^4x^4-\mathrm{C}_6^5x^5+\mathrm{C}_6^6x^6\,\, \left(*\right)
	\end{aligned}$\\
	Thay ${x=1}$ vào $(*)$, ta được $(1-1)^6=\mathrm{C}_6^0-\mathrm{C}_6^1+\mathrm{C}_6^2-\mathrm{C}_6^3+\mathrm{C}_6^4-\mathrm{C}_6^5+\mathrm{C}_6^6=S$. Vậy $S=0$.
}
\end{ex}

\begin{ex}%[0D8H3-5]
	Cho $\left(1-\dfrac{1}{2}x\right)^5=a_0+a_1x+a_2x^2+a_3x^3+a_4x^4+a_5x^5$.
	\choiceTF
	{$a_3=\dfrac{5}{2}$}
	{\True $a_5=-\dfrac{1}{32}$}
	{\True Hệ số lớn nhất trong tất cả hệ số là $\dfrac{5}{2}$}
	{Tổng $a_0+a_1+a_2+a_3+a_4+a_5=\dfrac{1}{16}$}
\loigiai{
	$\begin{aligned}
		\text{Ta có }
		\left(1-\dfrac{1}{2} x\right)^5
		&=\mathrm{C}_5^0+\mathrm{C}_5^1\left(-\dfrac{1}{2} x\right)+\mathrm{C}_5^2\left(-\dfrac{1}{2} x\right)^2+\mathrm{C}_5^3\left(-\dfrac{1}{2} x\right)^3+\mathrm{C}_5^4\left(-\dfrac{1}{2} x\right)^4+\mathrm{C}_5^5\left(-\dfrac{1}{2} x\right)^5 \\
		&=1-\dfrac{5}{2} x+\dfrac{5}{2} x^2-\dfrac{5}{4} x^3+\dfrac{5}{16} x^4-\dfrac{1}{32} x^5 \\
		&=a_0+a_1 x+a_2 x^2+a_3 x^3+a_4 x^4+a_5 x^5 \quad (*)
	\end{aligned}$\\
	Suy ra $a_0=1$, $a_1=-\dfrac{5}{2}$, $a_2=\dfrac{5}{2}$, $a_3=-\dfrac{5}{4}$, $a_4=\dfrac{5}{16}$, $a_5=-\dfrac{1}{32}$. \\
	Ta thấy hệ số lớn nhất tìm được là $a_2=\dfrac{5}{2}$. \\
	Thay $x=1$ vào $(*)$, ta được: $\left(1-\dfrac{1}{2}\right)^5=a_0+a_1+a_2+a_3+a_4+a_5$. \\
	Vậy $a_0+a_1+a_2+a_3+a_4+a_5=\dfrac{1}{32}$.
}
\end{ex}
\Closesolutionfile{ans}
\indapan{3}{ans/0C8B3-1}

\Opensolutionfile{ans}[ans/ans-0-B1-KQ]
\caukq
\begin{ex}%[0D8V3-4]
	Tìm hệ số của $x^5$ trong khai triển biểu thức $x(2x-1)^6+(x-3)^8$.
	\shortans{$-1272$}
	\loigiai{
		Hệ số của $x^5$ trong khai triển biểu thức $x(2x-1)^6$ là $\mathrm{C}_6^4 2^4(-1)^2=240$;\\
		Hệ số của $x^5$ trong khai triển biểu thức $(x-3)^8$ là $\mathrm{C}_8^5(-3)^3=-1512$.\\
		Suy ra hệ số của $x^5$ trong khai triển biểu thức $x(2x-1)^6+(x-3)^8$ là $240-1512=-1272$.
	}
\end{ex}
\begin{ex}%[0D8V3-6]
	Cho $n$ là số nguyên dương thỏa
	$\mathrm{C}_n^0+4\mathrm{C}_n^1+4^2\mathrm{C}_n^2+\ldots+4^n\mathrm{C}_n^n=15625$. Tìm $n$.
	\shortans{$6$}
	\loigiai{
		Xét khai triển $(1+x)^n=\mathrm{C}_n^0+\mathrm{C}_n^1x+\mathrm{C}_n^2x^2+\ldots+\mathrm{C}_n^nx^n$.\\
		Chọn $x=4$, ta có $5^n=\mathrm{C}_n^0+4 \mathrm{C}_n^1+4^2 \mathrm{C}_n^2+\ldots+4^n \mathrm{C}_n^n$.\\
		Suy ra $15625=5^n \Leftrightarrow 5^6=5^n \Leftrightarrow n=6$.
	}
\end{ex}
\begin{ex}%[0D8V3-6]
	Tìm số nguyên dương $n$ thỏa mãn $\mathrm{C}_{2n}^0+\mathrm{C}_{2n}^2+\mathrm{C}_{2n}^4+\mathrm{C}_{2n}^6+\cdots \cdots+\mathrm{C}_{2n}^{2n}=512$.
	\shortans{$5$}
	\loigiai{
		Xét $(1+1)^{2n}=\sum_{k=0}^{2n} \mathrm{C}_{2n}^k \cdot 1^{2n-k} \cdot 1^k=\mathrm{C}_{2n}^0+\mathrm{C}_{2n}^1+\mathrm{C}_{2n}^2+\mathrm{C}_{2n}^3+\mathrm{C}_{2n}^4+\cdots+\mathrm{C}_{2n}^{2n}$.\\
		Xét $(1-1)^{2n}=\sum_{k=0}^{2n} \mathrm{C}_{2n}^k \cdot 1^{2n-k} \cdot(-1)^k=\mathrm{C}_{2n}^0-\mathrm{C}_{2n}^1+\mathrm{C}_{2n}^2-\mathrm{C}_{2n}^3+\mathrm{C}_{2n}^4-\cdots+\mathrm{C}_{2n}^{2n}$.\\
		Lấy $(1)$ cộng $(2)$ ta được
		\begin{eqnarray*}
			&&2^{2n}+0^{2n}=2 \cdot\left(\mathrm{C}_{2n}^0+\mathrm{C}_{2n}^2+\mathrm{C}_{2n}^4+\mathrm{C}_{2n}^6+\cdots+\mathrm{C}_{2n}^{2n}\right) \\
			&\Leftrightarrow& 2^{2n}=2\cdot 512 \Leftrightarrow 2^{2n}=1024 \\
			&\Leftrightarrow& 2^{2n}=2^{10}\Leftrightarrow 2n=10\Leftrightarrow n=5.
		\end{eqnarray*}
	}
\end{ex}
\begin{ex}%[0D8V3-5]
	Cho tổng  $S=\mathrm{C}_{22}^{12}+\mathrm{C}_{22}^{13}+\ldots+\mathrm{C}_{22}^{20}+\mathrm{C}_{22}^{21}+\mathrm{C}_{22}^{22}$. Tìm ba chữ số đầu của $S$.
	\shortans{$174$}
	\loigiai{
		Ta có $2^{22}=(1+1)^{22}=\mathrm{C}_{22}^0+\mathrm{C}_{22}^1+\mathrm{C}_{22}^2+\ldots+\mathrm{C}_{22}^{20}+\mathrm{C}_{22}^{21}+\mathrm{C}_{22}^{22}$.\\
		Áp dụng $\mathrm{C}_n^k=\mathrm{C}_n^{n-k}$, ta được $\mathrm{C}_{22}^0=\mathrm{C}_{22}^{22}, \mathrm{C}_{22}^1=\mathrm{C}_{22}^{21},\ldots, \mathrm{C}_{22}^2=\mathrm{C}_{22}^{20}, \ldots, \mathrm{C}_{22}^{10}=\mathrm{C}_{22}^{12}$.\\
		Khi đó, ta có
		\begin{eqnarray*}
			&&\mathrm{C}_{22}^0+\mathrm{C}_{22}^1+\mathrm{C}_{22}^2+\ldots+\mathrm{C}_{22}^{20}+\mathrm{C}_{22}^{21}+\mathrm{C}_{22}^{22}=2\left(\mathrm{C}_{22}^{12}+\mathrm{C}_{22}^{13}+\ldots+\mathrm{C}_{22}^{20}+\mathrm{C}_{22}^{21}+\mathrm{C}_{22}^{22}\right)+\mathrm{C}_{22}^{11} \\
			&\Leftrightarrow& \mathrm{C}_{22}^{12}+\mathrm{C}_{22}^{13}+\ldots+\mathrm{C}_{22}^{20}+\mathrm{C}_{22}^{21}+\mathrm{C}_{22}^{22}=\dfrac{\mathrm{C}_{22}^0+\mathrm{C}_{22}^1+\mathrm{C}_{22}^2+\ldots+\mathrm{C}_{22}^{20}+\mathrm{C}_{22}^{21}+\mathrm{C}_{22}^{22}}{2}-\dfrac{\mathrm{C}_{22}^{11}}{2} \\
			&\Leftrightarrow& \mathrm{C}_{22}^{12}+\mathrm{C}_{22}^{13}+\ldots.+\mathrm{C}_{22}^{20}+\mathrm{C}_{22}^{21}+\mathrm{C}_{22}^{22}=\dfrac{2^{22}}{2}-\dfrac{\mathrm{C}_{22}^{11}}{2} \\
			&\Leftrightarrow& \mathrm{C}_{22}^{12}+\mathrm{C}_{22}^{13}+\ldots.+\mathrm{C}_{22}^{20}+\mathrm{C}_{22}^{21}+\mathrm{C}_{22}^{22}=2^{21}-\dfrac{\mathrm{C}_{22}^{11}}{2}\cdot
		\end{eqnarray*}
		Vậy $S=2^{21}-\dfrac{\mathrm{C}_{22}^{11}}{2}=1744436$. Ba chữ số đầu của $S$ là $174$.
	}
\end{ex}
\begin{ex}%[0D8V3-5]
	Cho biểu thức $S=3^{19}\mathrm{C}_{20}^0+3^{18}\mathrm{C}_{20}+3^{17}\mathrm{C}_{20}^2+\ldots+\dfrac{1}{3}\mathrm{C}_{20}^{20}$. Biết $S=\dfrac{a^b}{c}$, tính giá trị của $a\cdot b+c^2$.
	\shortans{$89$}
	\loigiai{
		$\begin{aligned}
			\text{Ta có }
			S&=3^{19}\mathrm{C}_{20}^0+3^{18}\mathrm{C}_{20}^1+3^{17}\mathrm{C}_{20}^2+\ldots +\dfrac{1}{3}\mathrm{C}_{20}^{20} \\
			\Leftrightarrow 3S&=3^{20}\mathrm{C}_{20}^0+3^{19}\mathrm{C}_{20}^1+3^{18}\mathrm{C}_{20}^2+\ldots +\mathrm{C}_{20}^{20}.
		\end{aligned}$ \\
		Xét khai triển \\
		$\begin{aligned}
			&(3+1)^{20}=\mathrm{C}_{20}^03^{20}1^0+\mathrm{C}_{20}^13^{19}1^1+\mathrm{C}_{20}^23^{18}1^2+\ldots +\mathrm{C}_{20}^{20}3^01^{20} \\
			\Rightarrow& (3+1)^{20}=\mathrm{C}_{20}^03^{20}+\mathrm{C}_{20}^13^{19}+\mathrm{C}_{20}^23^{18}+\ldots +\mathrm{C}_{20}^{20}\Rightarrow S=\dfrac{4^{20}}{3}.
		\end{aligned}$ \\
	Vậy $a\cdot b+c^2=89$.
	}
\end{ex}
\begin{ex}%[0D8C3-6]
	Một người có $500$ triệu đồng gửi tiết kiệm ngân hàng với lãi suất $7,2\%/$năm. Với giả thiết sau mỗi tháng người đó không rút tiền thì số tiền lãi được nhập vào số tiền ban đầu. Đây được gọi là hình thức lãi kép. Biết số tiền cả vốn lẫn lãi $T$ sau $n$ tháng được tính bởi công thức $T=T_0(1+r)^n$, trong đó $T_0$ là số tiênn gửi lúc đầu và $r$ là lãi suất của một tháng. Dùng hai số hạng đầu tiên trong khai triển của nhị thức Newtơn, tính gần đúng số tiền người đó nhận được (cả gốc lẫn lãi) sau $6$ tháng (đơn vị là triệu đồng).
	\shortans{$518$}
	\loigiai{
		Lãi suất của một tháng $r=\dfrac{7{,}2}{12} \%=0{,}6 \%/$tháng.\\
		$\begin{aligned}
			T=T_0(1+r)^n&=500\cdot 10^6(1+0{,}006)^6 \\
			&\approx 500\cdot 10^6\left(\mathrm{C}_6^0+\mathrm{C}_6^1 \cdot 0{,}006\right) \approx 518\, 000\, 000 \text{ đồng}.
		\end{aligned}$\\
		Vậy sau $6$ tháng người đó nhận được hơn $518$ triệu đồng.
	}
\end{ex}
\Closesolutionfile{ans}
\indapan{6}{ans/ans-0-B1-KQ}